\documentclass[sigconf,anonymous]{acmart}

% Remove ACM-specific metadata for draft
\settopmatter{printacmref=false}
\renewcommand\footnotetextcopyrightpermission[1]{}
\pagestyle{plain}

\usepackage{booktabs}
\usepackage{amsmath,amssymb}
\usepackage{algorithm}
\usepackage{algorithmic}
\usepackage{subcaption}
\usepackage{multirow}

\begin{document}

% ═══════════════════════════════════════════════════════════════
% TITLE
% ═══════════════════════════════════════════════════════════════

\title{Phantom: Distributional Forecasting for Cross-Sectional\\Cryptocurrency Return Prediction}

% Authors hidden for double-blind
\author{Anonymous}
\affiliation{\institution{Anonymous Institution}}

% ═══════════════════════════════════════════════════════════════
% ABSTRACT (~180 words)
% ═══════════════════════════════════════════════════════════════
%
% Structure: Problem → Gap → Approach → Results → Insight
% Must include: concrete quantitative OOS results
% Style: lead with finance problem, not ML technique

\begin{abstract}
Cross-sectional return prediction---ranking which assets will outperform their peers---is a cornerstone of quantitative investing, yet remains underexplored in cryptocurrency markets where distributional properties (heavy tails, regime shifts) demand more than point forecasts.
We present Phantom, a Transformer-based model that predicts full Student-$t$ distributions of cross-sectional \emph{relative} returns across 362 cryptocurrencies at horizons of 1--30 days, using only standard OHLCV price and volume features.
The model employs a patched encoder with cross-attention decoding for simultaneous multi-horizon prediction, trained on cross-sectionally demeaned returns to isolate idiosyncratic signal from shared market factors.
On 15 months of fully out-of-sample data (January 2025--March 2026), Phantom achieves a rank information coefficient (IC) of 0.124 at the 10-day horizon ($t = 8.93$, Newey-West), with a long-short portfolio Sharpe ratio of 11.84 net of 30 basis points transaction costs, a 74.7\% daily win rate, and perfect decile monotonicity.
The signal is orthogonal to the well-documented low-volatility anomaly.
Through systematic ablation across eight model versions, we establish that the cross-sectional signal resides in daily OHLCV patterns, cannot be extracted from absolute returns, does not benefit from derivative market features, and scales with the breadth of the asset universe---but not with higher-frequency data.
\end{abstract}

\maketitle

% ═══════════════════════════════════════════════════════════════
% 1. INTRODUCTION (~1.2 pages)
% ═══════════════════════════════════════════════════════════════
%
% Content plan:
% - Para 1: Crypto markets are large, volatile, 24/7. Cross-sectional
%   prediction (ranking) is more tractable than directional (absolute).
%   Cite Liu et al. 2022, Fieberg et al. 2024.
% - Para 2: Gap: existing work uses point forecasts and standard factors
%   (momentum, size, volume). No distributional cross-sectional model
%   exists. Why distributions matter: risk management, position sizing,
%   tail risk awareness.
% - Para 3: Our approach in 2-3 sentences. Transformer + Student-t +
%   relative returns + 362 assets.
% - Para 4: Key results (quantitative). IC, Sharpe, win rate, costs.
%   Emphasize OOS and statistical significance (NW t-stats).
% - Para 5: Ablation story. v1-v8 progression establishes what works
%   and what doesn't. Negative results as contributions.
% - Para 6: Contributions bullet list:
%   (1) First distributional model for cross-sectional crypto prediction
%   (2) Systematic ablation: synthetic pretraining, feature engineering,
%       granularity, universe breadth
%   (3) Strong OOS performance surviving transaction costs
%   (4) Orthogonality to known anomalies (low-vol)
% - Para 7: Paper organization.

\section{Introduction}

[TO FILL: Opening on crypto cross-sectional prediction opportunity.
Cite Liu et al. (2022) three-factor model, Fieberg et al. (2024) CTREND.
Contrast with equity literature. Emphasize 24/7 trading, high volatility,
market inefficiency relative to equities.]

[TO FILL: Gap --- no distributional cross-sectional models exist.
Point forecasts discard uncertainty information critical for position
sizing and risk management. Student-t captures heavy tails that
Gaussian models miss.]

[TO FILL: Our approach. Phantom: Transformer encoder-decoder,
Student-t output head, trained on relative returns across 362 crypto
assets. Multi-horizon simultaneous prediction (1-30 days).]

[TO FILL: Key results. IC=0.124 (NW t=8.93), Sharpe 11.84 net of
30 bps, win rate 74.7\%, perfect decile monotonicity. All fully OOS
on 15 months of data.]

[TO FILL: Ablation narrative. Eight versions systematically test:
synthetic vs real pretraining, absolute vs relative returns,
additional features, data granularity, universe breadth.]

\paragraph{Contributions.}
\begin{itemize}
    \item We propose Phantom, the first distributional forecasting model
    for cross-sectional cryptocurrency return prediction, producing
    calibrated Student-$t$ distributions at 30 horizons simultaneously.

    \item Through systematic ablation across eight model versions, we
    establish that (i) cross-sectional signal exists in daily OHLCV
    features, (ii) absolute returns are unpredictable from price/volume
    alone (consistent with weak-form efficiency), (iii) derivative market
    features (funding rates, taker buy ratios) provide no incremental
    signal, and (iv) higher-frequency (4-hour) data degrades performance
    due to sample overlap.

    \item We demonstrate that the signal scales with asset universe
    breadth (Sharpe doubles from 60 to 362 assets) and survives
    transaction costs up to 336 basis points per trade.

    \item We show orthogonality to the low-volatility anomaly via
    double-sorted portfolios, establishing that Phantom captures a
    distinct source of cross-sectional alpha.
\end{itemize}

% ═══════════════════════════════════════════════════════════════
% 2. RELATED WORK (~0.8 pages)
% ═══════════════════════════════════════════════════════════════
%
% Content plan:
% - 2.1 Cross-sectional crypto factors: Liu et al. 2022 (size, momentum,
%   volume), Fieberg et al. 2024 (CTREND from 28 technical indicators),
%   carry strategies (BIS WP 1087). Establish that momentum is the
%   dominant known factor.
% - 2.2 Time series forecasting with Transformers: PatchTST (Nie et al.
%   2023), iTransformer, Chronos (Ansari et al. 2024). Foundation models
%   for time series. JointFM (Hackmann 2026) for distributional prediction
%   on synthetic data.
% - 2.3 Probabilistic forecasting: DeepAR, TFT, normalizing flows.
%   Student-t vs Gaussian vs mixture heads. CRPS as proper scoring rule.
% - 2.4 Position our work: first to combine distributional output with
%   cross-sectional ranking for crypto. Different from JointFM (synthetic
%   only, no real-data cross-sectional signal).

\section{Related Work}

\subsection{Cross-Sectional Cryptocurrency Factors}
[TO FILL: Liu et al. 2022 --- CMKT, CSMB, CMOM factors. Momentum
dominant but declining. Fieberg et al. 2024 --- CTREND aggregates
28 technical signals. BIS WP 1087 --- carry from funding rates.
Low-vol anomaly in crypto.]

\subsection{Transformer-Based Time Series Forecasting}
[TO FILL: PatchTST patched input, channel-independence. iTransformer
inverts attention. Chronos tokenized forecasting. JointFM synthetic
SDE distributional prediction. Our architecture closest to PatchTST
but with cross-attention decoder for multi-horizon + Student-t head.]

\subsection{Probabilistic Forecasting}
[TO FILL: DeepAR, TFT, normalizing flows. Student-t vs Gaussian.
CRPS as proper scoring rule (Gneiting \& Raftery 2007). Energy distance.
Why parametric (Student-t) is better than quantile regression for
financial applications.]

% ═══════════════════════════════════════════════════════════════
% 3. METHODOLOGY (~2 pages)
% ═══════════════════════════════════════════════════════════════
%
% This is the technical core. Must be precise and mathematical.

\section{Methodology}

\subsection{Problem Formulation}
% Define:
% - Universe of N assets, daily OHLCV observations
% - Context window x_{t-L:t} of L=120 days, C=6 channels
% - Target: relative return r^{rel}_{i,t+h} = r_{i,t+h} - \bar{r}_{t+h}
%   where \bar{r}_{t+h} = (1/N_t) \sum_j r_{j,t+h}
% - Model outputs p(r^{rel}_{i,t+h} | x_{i,t-L:t}) = Student-t(\mu, \sigma, \nu)
%   for h = 1, ..., H (H=30)
% - Cross-sectional ranking: rank assets by \mu_{i,h} on each date

[TO FILL: Formal mathematical problem definition. Define notation
for assets $i \in \{1, \ldots, N_t\}$, context $\mathbf{x}_i \in \mathbb{R}^{L \times C}$,
relative return $r^{\text{rel}}_{i,h}$, predictive distribution
$p_\theta(r^{\text{rel}}_{i,h} | \mathbf{x}_i)$. State the dual objective:
(1) calibrated distributional forecasts, (2) cross-sectional ranking signal.]

\subsection{Feature Engineering}
% 6 OHLCV channels:
% Ch 0: log return = log(c_t / c_{t-1})
% Ch 1: intraday range = (h_t - l_t) / c_t  (Parkinson vol proxy)
% Ch 2: body ratio = (c_t - o_t) / (h_t - l_t)  (candle body)
% Ch 3: log volume ratio = log(v_t / median_{30}(v))
% Ch 4: trailing vol = std_{30}(r) * sqrt(252)
% Ch 5: momentum = sum_{10}(r)
%
% Emphasize: NO normalization (RevIN/MeanAbsScaling destroy signal)
% All features are asset-universal (work for any tradeable asset)

[TO FILL: Mathematical definition of each channel. Table of channels
with name, formula, intuition. Discuss why no normalization.]

\subsection{Model Architecture}
% Patched Transformer encoder:
%   - Input: x \in R^{L x C}, patch_len=5, stride=5 -> N=24 patches
%   - Patch embedding: flatten (C x patch_len) -> Linear -> d_model=512
%   - Positional encoding (learned)
%   - 8-layer Transformer encoder (pre-norm, GELU, d_ff=2048)
%
% Cross-attention decoder:
%   - Horizon embedding: h -> R^{d_model} (learned, h=1..30)
%   - 2-layer cross-attention: query=horizon_embed, key/value=encoder_out
%   - All 30 horizons decoded simultaneously (batched queries)
%
% Student-t head:
%   - Linear(d_model) -> (\mu, log\sigma, log\nu)
%   - sigma = exp(log_sigma), clamped to [exp(-7), exp(2)]
%   - nu = exp(log_nu), clamped to [2.01, 100]
%   - Output: Student-t(mu, sigma, nu) per horizon
%
% Anti-collapse mechanisms:
%   - Condition dropout (p=0.15): zero encoder output with probability p
%     during training -> classifier-free guidance at inference
%   - Encoder variance penalty: encourage diverse representations
%
% Architecture figure reference

[TO FILL: Full mathematical description of each component. Include
architecture diagram (Figure). Equation for patch embedding,
cross-attention, Student-t parameterization. Discuss anti-collapse.]

\subsection{Training Objective}
% Multi-component loss:
%   L = L_NLL + 0.5 * L_CRPS + 0.3 * L_MSE + 0.1 * L_var
%
% L_NLL = -log p(y | mu, sigma, nu)  (Student-t NLL)
% L_CRPS = closed-form CRPS for Student-t
% L_MSE = weighted MSE on predicted mean, sqrt(h) weighting, h >= 5
% L_var = encoder variance penalty
%
% Horizon weighting: sqrt(h/H) emphasizes longer horizons
%   where SNR ~ sqrt(h)

[TO FILL: Loss function equations. Explain each component and why
it's needed. CRPS definition and why it's a proper scoring rule.
Horizon weighting motivation from drift SNR scaling.]

\subsection{Cross-Sectional Demeaning}
% Relative return computation:
%   r^{rel}_{i,t,h} = r^{abs}_{i,t,h} - (1/N_t) * sum_j r^{abs}_{j,t,h}
%
% Why: removes shared market factor (beta), isolates idiosyncratic signal
% Computed offline over all (date, asset) pairs
% 99.9% of samples adjusted (groups with >= 3 assets)

[TO FILL: Mathematical formulation. Explain why this removes
market beta. Connection to Fama-French methodology. Discuss
minimum group size.]

% ═══════════════════════════════════════════════════════════════
% 4. EXPERIMENTAL SETUP (~0.8 pages)
% ═══════════════════════════════════════════════════════════════

\section{Experimental Setup}

\subsection{Data}
% 362 crypto assets from Binance (all active USDT pairs)
% Filter: remove stablecoins, leveraged tokens, < 300 days history
% Daily OHLCV, 2017-08 to 2026-03
% Split: train < 2024 (230K), val = 2024 (97K), test >= 2025 (145K)
% 1-year gap between train and test (val acts as buffer)
% Table: dataset statistics

[TO FILL: Data description table. Asset count, sample counts per
split, date ranges. Discuss Binance as data source. Filter criteria.]

\subsection{Baselines}
% 1. Momentum (5d, 10d, 20d): rank by past cumulative return
% 2. Mean-reversion (1d, 5d): rank by negative past return
% 3. Abnormal volume (5d): rank by log volume ratio
% 4. Low volatility: rank by negative trailing vol
% 5. Random: shuffled signal (null hypothesis)
%
% All evaluated with same L/S methodology (top/bottom 20%)

[TO FILL: Baseline descriptions. Cite momentum literature.
Explain why these are the natural comparisons.]

\subsection{Evaluation Metrics}
% Cross-sectional signal:
%   - Rank IC (Spearman) per date, averaged
%   - Newey-West t-statistic for IC and L/S returns
%   - Long-short portfolio: top/bottom quintile, equal-weighted
%   - Sharpe, Sortino, Calmar ratios (annualized with 365 days)
%   - Win rate, max drawdown
%
% Distributional quality:
%   - NLL, CRPS
%   - PIT histogram (should be uniform)
%   - Coverage calibration
%   - KS test for PIT uniformity
%
% Robustness:
%   - Monthly IC stability
%   - Rolling 60-day Sharpe
%   - Decile monotonicity (Spearman of decile returns with decile rank)
%   - Double-sort orthogonality vs low-vol
%   - Transaction cost analysis with turnover

[TO FILL: Mathematical definitions of each metric. Cite Gneiting \&
Raftery (2007) for CRPS, Newey \& West (1987) for t-stats.]

% ═══════════════════════════════════════════════════════════════
% 5. RESULTS (~2 pages)
% ═══════════════════════════════════════════════════════════════
%
% This is where all 4 figures and key tables go.

\section{Results}

\subsection{Cross-Sectional Signal}
% Table 1: IC by horizon (1d to 30d) with NW t-stats
% Figure 1a: IC by horizon bar chart
% Key finding: IC increases with horizon (0.08 at 1d to 0.17 at 30d)
%   consistent with drift SNR ~ sqrt(h)
% IC positive 78% of days, 15/15 months

[TO FILL: Present IC results. Table with horizons, IC, NW t-stat.
Discuss SNR scaling with horizon. Compare to crypto factor literature
(typical IC ~ 0.03-0.05 for known factors).]

\subsection{Portfolio Performance}
% Table 2: Strategy performance summary
%   Sharpe (gross/net), Sortino, Calmar, win rate, turnover, MDD
% Figure 1b: Cumulative L/S returns (gross, net 10bps, net 30bps)
% Figure 1c: Decile portfolio returns (perfect monotonicity)
% Quintile sort table with NW t-stats
%   Q5-Q1 spread = 2.53%/day, NW t = 7.44, Sharpe = 13.0
% Breakeven cost = 336 bps

[TO FILL: Present portfolio results. Quintile table in publication
format. Discuss why Sharpe is high (diversification across 362 assets,
Sharpe ~ IC * sqrt(N)). Transaction cost analysis table.]

\subsection{Baseline Comparison}
% Table 3: Model vs baselines (IC, Sharpe, win rate)
% Figure 1d: IC bar chart
% Key: Phantom IC=0.124 vs best baseline (low-vol IC=0.165)
% But: double-sort shows orthogonality (next subsection)
% Momentum has negative IC (mean-reverting market in test period)

[TO FILL: Baseline comparison table. Discuss why momentum is negative
(test period characteristics). Note that low-vol is strong.]

\subsection{Orthogonality to Low-Volatility Anomaly}
% Figure 4a: Double-sort heatmap (vol tercile x Phantom tercile)
% Figure 4b: Combined signal analysis
% Key: Phantom L/S spread significant within EACH vol tercile
%   -> signal is independent of low-vol
% Signal correlation between Phantom and low-vol
% Combined signal IC and Sharpe

[TO FILL: Double-sort results. Show that Phantom's spread has
significant NW t-stat within each vol tercile. Discuss signal
correlation. Combined signal analysis.]

\subsection{Distributional Quality}
% Figure 3: PIT histogram, coverage calibration, sigma term structure
% Coverage 90% -> 90.3% (well-calibrated)
% Sigma follows sqrt(h) scaling at short horizons
% Nu ~ 2: heavy-tailed Student-t (heavier than Gaussian)

[TO FILL: Distributional metrics. PIT uniformity, coverage,
sigma term structure. Discuss nu ~ 2 and heavy tails.]

\subsection{Robustness}
% Figure 2: Monthly IC, rolling Sharpe, turnover
% IC positive all 15 months
% Rolling 60-day Sharpe always > 0 (min = 2.82)
% Turnover stable at 38%

[TO FILL: Robustness checks. Monthly stability, rolling Sharpe,
turnover analysis.]

\subsection{Ablation: What Works and What Doesn't}
% Table 4: Version progression (v1-v8)
%   v1/v2: synthetic pretraining -> zero transfer
%   v3: real multi-asset -> calibration, no directional signal
%   v4: absolute returns -> memorization, doesn't generalize
%   v5: relative returns -> IC=0.09 (breakthrough)
%   v6: + new features -> features don't help, crypto-only helps
%   v7: 4h bars -> worse (sample overlap)
%   v8: 362 assets -> Sharpe doubles
%
% Key lessons as bullet points

[TO FILL: Ablation table. Each version: what changed, IC, Sharpe.
Discuss each negative result as a contribution. Emphasize that
the progression is a systematic search over the design space.]

% ═══════════════════════════════════════════════════════════════
% 6. CONCLUSION (~0.3 pages)
% ═══════════════════════════════════════════════════════════════

\section{Conclusion}

[TO FILL:
- Summary: distributional Transformer for cross-sectional crypto,
  IC=0.124, Sharpe 11.84 net, orthogonal to low-vol
- Key insight: relative returns unlock signal that absolute returns
  cannot access; signal scales with universe breadth
- Limitations: crypto-only (equity IC~0), no conditional signal
  beyond ranking, no live trading validation
- Future work: live deployment, alternative data (on-chain),
  extension to DeFi protocols, adaptive position sizing using
  the distributional output]

% ═══════════════════════════════════════════════════════════════
% REFERENCES
% ═══════════════════════════════════════════════════════════════

\bibliographystyle{ACM-Reference-Format}

\begin{thebibliography}{99}

\bibitem{liu2022}
Y.~Liu, A.~Tsyvinski, and X.~Wu.
\newblock Common risk factors in cryptocurrency.
\newblock {\em Journal of Finance}, 77(2):1133--1177, 2022.

\bibitem{fieberg2024}
C.~Fieberg, T.~Hornuf, G.~Lau, and A.~Mihailova.
\newblock A trend factor for the cross section of cryptocurrency returns.
\newblock {\em Journal of Financial and Quantitative Analysis}, 2024.

\bibitem{nie2023}
Y.~Nie, N.~H. Nguyen, P.~Sinthong, and J.~Kalagnanam.
\newblock A time series is worth 64 words: Long-term forecasting with transformers.
\newblock {\em ICLR}, 2023.

\bibitem{ansari2024}
A.~F. Ansari, L.~Stella, et~al.
\newblock Chronos: Learning the language of time series.
\newblock {\em arXiv:2403.07815}, 2024.

\bibitem{hackmann2026}
D.~Hackmann.
\newblock JointFM: A joint foundation model for distributional prediction.
\newblock {\em arXiv:2603.20266}, 2026.

\bibitem{gneiting2007}
T.~Gneiting and A.~E. Raftery.
\newblock Strictly proper scoring rules, prediction, and estimation.
\newblock {\em Journal of the American Statistical Association}, 102(477):359--378, 2007.

\bibitem{newey1987}
W.~K. Newey and K.~D. West.
\newblock A simple, positive semi-definite, heteroskedasticity and autocorrelation consistent covariance matrix.
\newblock {\em Econometrica}, 55(3):703--708, 1987.

\bibitem{szekely2013}
G.~J. Sz{\'e}kely and M.~L. Rizzo.
\newblock Energy statistics: A class of statistics based on distances.
\newblock {\em Journal of Statistical Planning and Inference}, 143(8):1249--1272, 2013.

\bibitem{kim2022}
T.~Kim, J.~Kim, Y.~Tae, C.~Park, J.~Choi, and J.~Choo.
\newblock Reversible instance normalization for accurate time-series forecasting against distribution shift.
\newblock {\em ICLR}, 2022.

\bibitem{bis2024}
I.~Aldasoro, S.~Barth, and M.~Tsomocos.
\newblock Crypto carry.
\newblock {\em BIS Working Paper 1087}, 2024.

\end{thebibliography}

\end{document}
